\section{Experimental Results Summary}
\label{sec:results}

% Load the automatically compiled results table
% Generated by compile_paper_results.py. Do not edit manually.
\begin{table*}[t]
\centering
\caption{Performance Comparison and Discovery Metrics of SBI-Factory across Synthetic and Real-World domains.}
\label{tab:results}
\begin{tabular}{llccccr}
\toprule
Domain & Type & Starting MSE & Best Final MSE & Mean Final MSE & Imp. Runs & Accepted Mut.\\
\midrule
\multicolumn{7}{l}{\textbf{Synthetic Domains}} \\
Ecology & Synthetic & $16.5498$ & $2.7507$ & $6.5746 \pm 1.7906$ & 10/10 & 37 \\
Oncology & Synthetic & $0.0451$ & $0.0451$ & $0.0451 \pm 0.0000$ & 9/10 & 11 \\
Econ & Synthetic & $0.0318$ & $0.0315$ & $0.0318 \pm 0.0002$ & 0/10 & 0 \\
Battery & Synthetic & $0.4187$ & $0.0312$ & $0.0446 \pm 0.0202$ & 10/10 & 87 \\
PK/PD & Synthetic & $0.0255$ & $0.0255$ & $0.0255 \pm 0.0000$ & 0/10 & 0 \\
Synbio & Synthetic & $0.1977$ & $0.0672$ & $0.1217 \pm 0.0469$ & 9/10 & 18 \\
Climate & Synthetic & $0.2602$ & $0.1622$ & $0.2416 \pm 0.0366$ & 2/10 & 4 \\
Cardio & Synthetic & $11.7027$ & $11.1246$ & $11.7027 \pm 0.3448$ & 0/10 & 0 \\
Belousov-Zhabotinsky & Synthetic & $\infty$ & $0.4399$ & $0.7962 \pm 0.2420$ & 10/10 & 72 \\
Epidemiology & Synthetic & $0.8466$ & $0.8455$ & $0.8457 \pm 0.0001$ & 10/10 & 82 \\
Neuro & Synthetic & $0.6372$ & $0.6311$ & $0.6366 \pm 0.0018$ & 9/10 & 28 \\
Fluid & Synthetic & $\infty$ & $59.1984$ & $519.8883 \pm 320.6920$ & 10/10 & 97 \\
Astro & Synthetic & $0.3313$ & $0.2932$ & $0.3206 \pm 0.0162$ & 10/10 & 35 \\
Chem Kinetics & Synthetic & $0.2509$ & $0.1049$ & $0.2074 \pm 0.0664$ & 10/10 & 51 \\
Immune & Synthetic & $3.2986$ & $0.0328$ & $0.5048 \pm 0.5850$ & 10/10 & 47 \\
Pop Genetics & Synthetic & $0.4317$ & $0.3889$ & $0.4264 \pm 0.0125$ & 10/10 & 108 \\
Thermo & Synthetic & $0.8708$ & $0.0020$ & $0.0020 \pm 0.0000$ & 10/10 & 65 \\
Materials & Synthetic & $1.0429$ & $0.1572$ & $0.8035 \pm 0.3252$ & 10/10 & 55 \\
Agriculture & Synthetic & $0.0255$ & $0.0233$ & $0.0251 \pm 0.0007$ & 9/10 & 37 \\
Social & Synthetic & $0.1051$ & $0.0215$ & $0.0391 \pm 0.0257$ & 10/10 & 144 \\
\midrule
\multicolumn{7}{l}{\textbf{Real-World Domains}} \\
Solar Sunspots & Real-World & $19.3126$ & $1.4903$ & $5.7296 \pm 4.0442$ & 10/10 & 75 \\
Nile River Volume & Real-World & $1678.6353$ & $0.4204$ & $24.7047 \pm 26.2223$ & 10/10 & 65 \\
US Real GDP & Real-World & $336.1619$ & $12.6558$ & $27.4079 \pm 11.6833$ & 10/10 & 78 \\
Oral Theophylline & Real-World & $10.5120$ & $0.0675$ & $0.0865 \pm 0.0086$ & 10/10 & 87 \\
Atmospheric $\text{CO}_2$ & Real-World & $1.1230$ & $0.0069$ & $0.2634 \pm 0.1908$ & 10/10 & 68 \\
\bottomrule
\end{tabular}
\end{table*}


As shown in Table~\ref{tab:results}, the $E^3$ system achieved significant error reductions in many domains, while several domains showed negligible movement despite multiple accepted mutations. The table should therefore be read as a mixed-domain audit rather than a uniformly successful benchmark.

To evaluate $E^3$ against established symbolic regression and system identification methods, we maintain a preliminary baseline runner for \textbf{PySR} (numerical gradient approximation followed by symbolic search and trajectory rollout) and \textbf{PySINDy} (sparse identification of non-linear dynamics). The runner now saves train/test trajectory metrics and derivative-fit diagnostics for representative continuous domains. These baselines are not yet a final method-to-method benchmark because the $E^3$ historical artifacts were optimized with in-sample ABC-SMC distances, while the baseline runner reports held-out trajectory rollouts.

\textbf{Comparative Analysis Status:} 
\begin{itemize}
    \item \textbf{PySINDy} is evaluated by fitting on the training prefix and rolling the discovered system over both the train prefix and held-out suffix.
    \item \textbf{PySR} is fit to local numerical derivatives, then its learned right-hand-side predictors are integrated with \texttt{solve\_ivp} to produce trajectory-level train/test errors. Derivative-local errors are saved only as diagnostics.
    \item \textbf{$E^3$} should be compared against those baselines only after rerunning selected domains with matched train/test splits, fixed seeds, saved prompts, and compute-normalized budgets.
\end{itemize}

\textbf{Scope of Baselines:} Direct method-to-method comparisons using PySR and PySINDy are currently restricted to representative continuous domains. The present baseline script is suitable for saved preliminary trajectory diagnostics, but final claims require regenerated $E^3$ runs under the same held-out evaluation protocol.

It is critical to normalize for underlying compute: each structural evaluation in the Qwen runs used 15 generations of ABC-SMC with 150,000 initial particles and resampled subsequent generations. While $E^3$ reduces the need for blind symbolic tree mutation, its simulator evaluation footprint remains large, necessitating compute-normalized benchmarking in future work.

\subsection{Limitations and Failure Modes}
While $E^3$ successfully discovered underlying physical invariants across diverse domains, the system exhibits several fundamental failure modes that must be addressed in future work:
\begin{enumerate}
    \item \textbf{Stiff ODE Hallucinations:} LLMs occasionally propose mathematical structures containing extreme exponential terms or divergent asymptotic limits. When injected into the JAX/Diffrax pipeline, these equations create infinitely stiff systems that cause the adaptive Tsit5 solver to hang or exhaust memory evaluating vanishing step sizes. While we mitigate this by setting explicit \texttt{max\_steps} constraints, extreme hallucinations still result in wasted compute cycles.
    \item \textbf{Curse of Dimensionality in ABC-SMC:} While Covariate-Proportional regularization prevents early collapse, scaling the state vector beyond 50 coupled interacting dimensions causes the ABC-SMC acceptance rate to plummet exponentially. Pure SMC struggles to sample effectively in highly parameterized neural spaces, limiting $E^3$ to interpretable, low-to-medium parameter regimes.
    \item \textbf{Reward Hacking:} In certain synthetic datasets without sufficient noise, the LLM discovered that it could drastically reduce distance not by finding the physical invariant, but by overfitting a highly parameterized high-frequency sine wave to interpolate between data points. This ``reward hacking'' motivates explicit parsimony penalties and held-out forecasting metrics in future runners.
\end{enumerate}

\section{Distributed Deployment as a Follow-up Paradigm}
\label{sec:future}

While a single machine running $E^3$ is novel and sufficient for multi-generational scientific discovery, we propose a \textbf{distributed, asynchronous follow-up architecture} for large-scale deployments targeting previously unsolved scientific datasets (e.g., full-scale genomic transcription networks).

In this volunteer-grid paradigm, a central server would act as a coordinator, hosting target datasets and serving the global best equations. Workers could pull tasks, execute local ABC-SMC iterations, and upload accepted mutations once safe execution and provenance controls exist.

\subsection{Security and Sandboxing}
Because the LLM generates raw Python/JAX code that is executed locally, a distributed system must prevent malicious actors from uploading remote code execution payloads. Future iterations of $E^3$ will execute worker-submitted models in strictly isolated WebAssembly (Wasm) or gRPC sandboxes, guaranteeing that only pure mathematical vector operations are passed to the host JAX compiler.

\section{Conclusion}
\label{sec:conclusion}

We presented $E^3$, a multi-generational framework that leverages the structural coding capabilities of LLMs and ABC-SMC-style simulation scoring to automate parts of scientific model search. The current implementation demonstrates a practical continuous-dynamics discovery loop across 25 domains, with strong results in some settings and clear failure modes in others. Future work should rerun selected domains with matched train/test evaluation, stronger baselines, safer generated-code execution, and explicit stochastic simulators before making broader claims.
